\documentclass[11pt,a4paper]{article}
\usepackage[margin=25mm]{geometry}
\usepackage[T1]{fontenc}
\usepackage[utf8]{inputenc}
\usepackage[british]{babel}
\usepackage{lmodern,microtype,amsmath,amssymb,amsthm,booktabs,tabularx,graphicx}
\usepackage{xcolor,caption,placeins,tikz}
\usetikzlibrary{arrows.meta,positioning,calc}
\usepackage[colorlinks=true,linkcolor=teal,citecolor=teal,urlcolor=teal]{hyperref}
\captionsetup{font=small,labelfont=bf}
\setlength{\parskip}{4pt}
\setlength{\parindent}{0pt}
\setlength{\emergencystretch}{2em}
\definecolor{ink}{HTML}{204B57}
\definecolor{pale}{HTML}{EDF4F4}
\definecolor{warm}{HTML}{B26338}
\tikzset{box/.style={draw=ink,rounded corners=3pt,fill=pale,align=center,inner sep=7pt},
 arrow/.style={-{Stealth},thick,draw=ink}}
\newtheorem{proposition}{Proposition}
\newcommand{\Maj}{\operatorname{MAJ}_3}
\newcommand{\F}{\mathbb F}
\usepackage{tocloft}
\setlength{\cftsecnumwidth}{2.6em}
\setlength{\cftsubsecnumwidth}{3.6em}
\renewcommand{\thesection}{S\arabic{section}}
\renewcommand{\thesubsection}{\thesection.\arabic{subsection}}
\renewcommand{\theequation}{S\arabic{equation}}
\hypersetup{pdftitle={Supplementary Information: Solutions to the challenges proposed by 0xPARC},pdfauthor={Alberto Hernandez-Espinosa}}
\title{Supplementary Information\\[5pt]
\large Solutions to the challenges proposed by 0xPARC}
\author{Alberto Hernandez-Espinosa}
\date{15 September 2026}
\begin{document}
\maketitle
\tableofcontents
\newpage
\section{Questions, conventions, and contribution}
We number the source's puzzles Q1--Q8 in order: list recovery, majority,
Fourier, modular power sums, range, exclusion of one, 64-bit factorisation,
and 4096-bit factorisation~\cite{challenge}. The Fourier text gives
32,768 complex slots but 65,536 in its rotation example; we address both.
The prime in Q4 is $q=2^{127}-1$. The distinct constraint field in Q5--Q8 is
$\F_P$, with
\begin{align*}
P={}&21888242871839275222246405745257275088548364400416034343698204186575808495617.
\end{align*}
External field values must be canonical integers in $[0,P-1]$.

This supplement states the exact conventions and gives the full arguments
behind the main text. The Boolean reconstruction is described in Section S4;
the finite verification record is in Section S10.

\section{Q1: recovering an integer list}
Let $v=(v_0,\ldots,v_{n-1})$ be positive integers and let an oracle return exact
dot products. With no bound, first query $(1,\ldots,1)$ to obtain $S=\sum_i v_i$.
For $n>1$, put $B=S+1$ and query $(1,B,\ldots,B^{n-1})$, obtaining
\[
 R=\sum_{i=0}^{n-1}v_iB^i.
\]
Since $0<v_i<B$, repeated division by $B$ recovers the digits without carries
between positions. A known bound $v_i\le U$ permits $B=U+1$ in advance and
one query. For $n=1$, the first query alone recovers the list.

\begin{proposition}
For $n\ge2$ with unbounded positive integers and rational query coefficients,
two queries suffice and one query cannot always suffice.
\end{proposition}
\begin{proof}
Sufficiency follows from the adaptive construction above. For any fixed
rational query vector $x$, clearing denominators leaves an integer linear map
from $\mathbb Z^n$ to $\mathbb Z$. Its rational nullspace has a nonzero integer
vector $z$. Choose positive $v_i$ sufficiently large that $v_i+z_i>0$.
Then $v$ and $v+z$ are different positive lists with the same answer.
\end{proof}

An idealized one-query solution uses coefficients linearly independent over
$\mathbb Q$. For example, $1,t,\ldots,t^{n-1}$ for a transcendental real $t$
are independent: equality of two answers would give a nonzero integer
polynomial vanishing at $t$. This is an injectivity argument in an exact-real
model, not a floating-point recovery algorithm. A bounded integer encoding
already needs roughly $n\log_2 B$ answer bits; reducing the number of queries
does not remove that information requirement.

\section{Q2: majority using only three-input majority}
Let $f$ be a monotone Boolean function and choose distinct active variables
$a,b,c$. Write $f_{b\gets a}$ for the function formed by substituting $a$ for
$b$, and similarly for the other identifications. Then
\[
 f=\Maj(f_{b\gets a},f_{c\gets b},f_{a\gets c}).
\]
If the three bits agree, every substitution leaves them unchanged. Otherwise
one substitution leaves the assignment unchanged, one increases a bit, and
one decreases a bit. By monotonicity the three outputs are $f$, a value at
least $f$, and a value at most $f$. Their majority is $f$.

Odd-input majority is monotone and self-dual: complementing all input bits
complements its output. Identification preserves both properties. A monotone
self-dual Boolean function on at most two active variables is a projection:
its values at $00,11$ are $0,1$, and the values at $01,10$ are opposite.
Thus recursion terminates in input wires and uses only $\Maj$ gates.

The construction represents subproblems by a positive odd total of weights
over the original coordinates. Identification merges weights, a weight greater
than half the total immediately returns its input, and otherwise the first
three active variables are used. Memoisation identifies equal weight tuples.
Every operand is an external input or an earlier gate; ordered duplicate ports
and free fan-out are permitted. Constants and inversions are absent.

With $T(2)=0$ and $T(k)\le1+3T(k-1)$, a loose unshared bound is
$T(n)\le(3^{n-2}-1)/2$ for $n\ge3$. This proves the requested existence for
2025 inputs but is not a compactness guarantee. No 2025-input circuit is
materialised here. Gate counts below refer only to emitted circuits. Exhaustive
checks cover odd sizes through 11; exact decision diagrams cover 13 and 15.
The symbolic target is built independently by a threshold recurrence.

For comparison, a naive majority of disjoint group majorities fails: the
groups $110,110,000$ produce local votes $1,1,0$ and hence output one, although
only four of the nine original inputs are one.

\section{Index deconvolution: representation, recovery and controls}
\subsection{Exact index sets and schema expansion}
Let $f:\{0,1\}^n\to\{0,1\}$ be a deterministic output column. States are
ordered by $\iota(x)=\sum_{i=0}^{n-1}2^ix_i$, with bit zero varying fastest.
The active index set $\mathcal I_f$ consists of all indices with output one.
Membership in this set specifies the complete output column.

For a schema $\sigma\in\{0,1,*\}^n$, let $J_1(\sigma)$ be its fixed-one
positions and $J_*(\sigma)$ its don't-care positions. Define the decimal
anchor and the offset family by
\[
 a_\sigma=\sum_{i\in J_1(\sigma)}2^i,\qquad
 \Omega_\sigma=\left\{\sum_{i\in J_*(\sigma)}2^i z_i:z_i\in\{0,1\}\right\}.
\]
The sumandos are the fillings of the schema's own don't-care positions;
$\Omega_\sigma$ is the numerical offset family these fillings generate.
The indices represented by the schema are exactly
$a_\sigma+\Omega_\sigma$. Fixed-one and free positions are disjoint, so
adding their encodings cannot create a carry. A covering family $\mathcal C$
represents the output precisely when
\[
\mathcal I_f=\bigcup_{\sigma\in\mathcal C}(a_\sigma+\Omega_\sigma).
\]
This is a union, since schemata may overlap. One must retain each anchor's
associated offset family; arbitrary cross-pairing of anchors and offsets
need not preserve the function. A cover by complete assignments always
exists, but a compact cover or a minimal cover is not guaranteed.

For three-input majority, $\mathcal C=\{11*,1*1,*11\}$ gives anchor/offset
pairs $(3,\{0,4\})$, $(5,\{0,2\})$, $(6,\{0,1\})$. The expansion is
$\{3,5,6,7\}$. All three inputs affect the function, although each schema
has a free position. Disconnected coordinates are free in every schema;
freedom in one schema does not imply disconnection.

\subsection{Recovery of functional inputs}
From the complete column, recover
\[
I_c(f)=\{i:\exists x,\ f(x)\ne f(x\oplus e_i)\},
\]
where $x\oplus e_i$ flips bit $i$. These are the functional inputs. If
$i\notin I_c(f)$, changing only bit $i$ never changes the output. Any two
states agreeing on $I_c(f)$ can be joined by a sequence of flips outside
$I_c(f)$, so their outputs agree. Therefore there is a well-defined reduced
function $g$ with $f(x)=g(x|_{I_c(f)})$. Its truth table is obtained by
restricting the full column and checking consistency across all free
assignments. Expanding it to the original coordinates reconstructs $f$
exactly. This proves functional recovery from a complete deterministic
repertoire. Gate identification can match $g$ to a declared family; multiple
names or parameterisations may describe the same function.

Testing all coordinate pairs costs at most $n2^{n-1}$ comparisons per
output column, with $O(2^n)$ storage when the full column is retained.
The method does not remove the exponential cost of supplying an arbitrary
complete repertoire. A compact schema family is a representation benefit,
not a universal guarantee that this inverse computation is inexpensive.

For a network with observed outputs for every input state, apply the same
argument to each output column to recover a functionally equivalent map.
Structurally declared but functionally inactive inputs are not identifiable.
If only a circuit's final output is observed, internal gates and their
wiring are not identifiable either. Partial or noisy observations do not
satisfy this exact-recovery premise.

\subsection{Majority checks and negative controls}
For odd $n$, setting exactly $(n-1)/2$ of the other bits to one shows that
every input of majority is functional. Complementation gives exactly
$2^{n-1}$ active indices. The added manuscript check evaluates the emitted
circuits for $n=1,3,5,7,9,11$, applies functional-input detection and column
reduction, and compares the result with the independent condition
$\sum_i x_i>n/2$ on every state. This covers $2+8+32+128+512+2048=2730$
states. The check also expands the three schemata above exactly.

A free-input control embeds three-input majority in five coordinates. All
32 states recover $I_c=\{0,1,2\}$ and the eight-entry majority table;
coordinates 3 and 4 do not affect the output. An incorrect-circuit control
uses a majority of three disjoint three-input majorities. It has nine
functional inputs and 256 active indices, yet its full output column differs
from nine-input majority. Thus neither the number of functional inputs nor
the size of the active set establishes equality; exact membership is checked.
The main text's $110,110,000$ assignment is one explicit disagreement.

\subsection{The five-input network and pattern figures}
The worked network uses the emitted five-input circuit, with
\begin{align*}
g_0&=\Maj(x_0,x_2,x_4),&g_1&=\Maj(x_0,g_0,x_3),\\
g_2&=\Maj(x_1,x_3,x_4),&g_3&=\Maj(x_0,x_1,g_2),\\
g_4&=\Maj(x_2,x_3,x_4),&g_5&=\Maj(x_1,x_2,g_4),\\
y&=g_6=\Maj(g_1,g_3,g_5).
\end{align*}
The figures show this circuit's connections, not a generic illustrative
network. All seven internal outputs are evaluated for all 32 inputs.
The output column is compared state by state with $\sum_i x_i\ge3$;
functional-input detection recovers all five coordinates.

To generate the schema family, choose every three-element subset of the
five input positions, fix those bits to one and leave the other two free.
This gives $\binom53=10$ schemata. Expanding their associated anchors and
offsets and taking the union gives precisely the circuit's 16 active indices.
The displayed $111**$ schema has anchor 7 and offsets $0,8,16,24$.
Its four complete input states, together with their internal outputs, are
\begin{center}
\begin{tabular}{rcc}\toprule
Index & Input $x_0\ldots x_4$ & Internal outputs $g_0\ldots g_6$\\\midrule
7 & 11100 & 1101011\\
15 & 11110 & 1111111\\
23 & 11101 & 1111111\\
31 & 11111 & 1111111\\\bottomrule
\end{tabular}
\end{center}
For each input, the figure-generation check retains a pair of states
that differ only at that input and have different outputs. The displayed
pair is index 6 ($01100$, output zero) and index 7 ($11100$, output one).
The complete repertoire, internal signals, schemata and checks accompany
the figures as numerical data. Both diagram topology and table entries
are derived from the same circuit, and every schema is checked for exact
membership before rendering.
These controls establish the stated finite reconstruction behaviour and
illustrate its limitations; they are not a new synthesis algorithm or a
claim to recover hidden internal topology.

\section{Q3: depth-three packed Fourier circuits}
We use the forward unnormalised DFT and rotation convention
\[
 (Fx)_r=\sum_{j=0}^{N-1}x_j e^{-2\pi i rj/N},\qquad
 R_k(x)[r]=x[(r+k)\bmod N].
\]
The primitives are vector addition, rotation by any offset, and componentwise
multiplication by a known complex vector. Their charged costs are respectively
4, 6100, and 5600 microseconds. Every known-vector multiplication consumes one
multiplication level, including masks. This explicitly states how constants
are charged under the puzzle's simplified model.

\subsection{A complete factorisation}
At DIF stage $s=0,\ldots,L-1$, where $N=2^L$, take blocks of length $m=N/2^s$
and $h=m/2$. For each block and $0\le t<h$, the stage maps
\[
 y_t=x_t+x_{t+h},\qquad
 y_{t+h}=(x_t-x_{t+h})e^{-2\pi i t/m}.
\]
The first output half contains the even-indexed frequencies after recursively
transforming its sum inputs; the second contains the odd-indexed frequencies
after transforming its twiddled differences. Induction on $L$ gives the DFT
in bit-reversed frequency order. Applying the bit-reversal permutation to the
output therefore yields the natural-order DFT.

Partition the $L$ stages into at most three nonempty contiguous intervals.
Each interval is fused into a linear transform; the final transform includes
the bit-reversal permutation. For interval $[a,b)$, let $m=N/2^a$ and $h=N/2^b$.
Before final permutation, each row $r$ has precisely these input columns:
\[
 \left\lfloor r/m\right\rfloor m+(r\bmod h)+ht,
 \qquad 0\le t<2^{b-a}.
\]
Following a row backwards through the butterfly stages gives one path to
each such column. An upper output contributes coefficients $+1,+1$; a lower
output contributes the stage twiddle times $+1,-1$. Products along the chosen
path determine the coefficient exactly. In the last transform replace $r$
inside this rule by its $L$-bit reversal. This gives explicit, lazily generated
coefficient vectors, without a dense full-size matrix.

\subsection{Charged diagonals and rotation reuse}
For each fused matrix $A$, set $d_k[r]=A[r,(r+k)\bmod N]$. With $K$ the set
of nonzero cyclic diagonals,
\[
 Ax=\sum_{k\in K}d_k\odot R_k(x).
\]
For a baby-step size $b_0\mid N$, write $k=gb_0+i$, $0\le i<b_0$, and regroup:
\[
 Ax=\sum_g R_{gb_0}\left(
       \sum_{i:\,gb_0+i\in K}R_{-gb_0}(d_{gb_0+i})\odot R_i(x)
                         \right).
\]
This identity follows by evaluating both sides at a slot: the outer rotation
cancels the opposite rotation of the known coefficient. Rotating known
coefficient data is part of its offline definition; every rotation of the
encrypted vector is an explicit charged node. If $I$ and $G$ are the used
baby and giant indices, the emitted transform contains $|K|$ multiplications,
$|K|-1$ additions, and $|I\setminus\{0\}|+|G\setminus\{0\}|$ rotations.
There is one multiplication level per fused transform.

For nonfinal intervals, $K$ is the set of $ht\bmod N$ with
$-(2^{b-a}-1)\le t\le2^{b-a}-1$. For the final interval, $h=1$, and each
row contributes a cyclic interval of $m$ offsets starting at
$\lfloor\operatorname{rev}_L(r)/m\rfloor m-r$. A cyclic interval union finds
the exact support without enumerating an $N\times N$ matrix.

\subsection{Search family and comparison}
All ordered contiguous partitions into one, two, or three groups are examined:
$1+(L-1)+\binom{L-1}{2}$, giving 106 and 121. For every block we evaluate every
power-of-two baby size dividing $N$. Costs add across blocks, so selecting the
best baby size independently per block evaluates the complete Cartesian
family without materializing redundant tuples. Ties use multiplication count,
rotation count, partition order, then baby-size order. The ledger retains every
evaluated block choice and its status. The result is best within this family,
not a globally optimal arithmetic circuit.

The dense depth-one baseline uses all $N$ diagonals with the same rotation
regrouping and cost model. Its best baby sizes are 128 and 256 for the two
dimensions (the symmetric 256 alternative ties at $N=32768$). Numerical
validation executes every emitted primitive, compares with an independent FFT,
and requires normalised maximum error at most $10^{-10}$. Every basis vector
through $N=128$ is also compared with the direct exponential formula. No
ciphertext latency is inferred from the plaintext runtime.

\section{Q4: modular power sums}
For the stated prime $q=2^{127}-1$, only the zero solution exists. The linear
equation gives $d=-(a+b+c)$. Expanding the cube equation gives
\[
 a^3+b^3+c^3+d^3=-3(a+b)(a+c)(b+c).
\]
Since $q\ne3$ and $\F_q$ has no zero divisors, one of the three pair sums is
zero. Relabel so $b=-a$ and hence $d=-c$. The square equation becomes
$2(a^2+c^2)=0$. Here $q\ne2$ and $q\equiv3\pmod4$. If $c\ne0$, then
$z=a/c$ would satisfy $z^2=-1$. But Fermat's theorem gives
$1=z^{q-1}=(-1)^{(q-1)/2}=-1$, a contradiction. Thus $c=0$ and then
$a=b=d=0$. This proof uses the primality stated in the puzzle.

Finite checks verify the polynomial identity over the integers and exhaustively
enumerates small fields. Primes 7, 11, and 19 have only zero. Characteristic
3 is exceptional: $(1,1,1,0)$ satisfies all three equations. For 5 and 13,
$-1$ has square roots 2 and 5 respectively, giving nonzero opposite-pair
solutions. These examples illustrate the assumptions; they do not substitute
for the proof at $q$.

\section{Q5 and Q6: range and exclusion constraints}
A quadratic row means $A(w)B(w)=C(w)$ in $\F_P$, with all three expressions
linear in the signals. For a 64-bit range use auxiliary bits $b_0,\ldots,b_{63}$:
\[
 b_i(b_i-1)=0\ (0\le i<64),\qquad x=\sum_{i=0}^{63}2^ib_i.
\]
Every satisfying bit is zero or one because a field has no zero divisors.
The sum lies in $[0,2^{64}-1]\subset[0,P-1]$, so its canonical representative
equals $x$ as an integer. Conversely the ordinary binary expansion of any
such $x$ provides a witness. A field circuit cannot distinguish $x$ from an
external integer $x+P$; canonical input validation is therefore explicit.

To exclude one, impose $(r-1)s=1$. It has a witness exactly when $r\ne1$:
take the inverse $s=(r-1)^{-1}$. At $r=1$ the equation would be $0=1$.

\section{Q7: a 64-bit factorisation}
Range-constrain public $n$ and private $u,v$ to 64 bits, enforce $uv=n$, and
enforce $(u-1)s=1$ and $(v-1)t=1$. The system also explicitly enforces
$u,v\ge2$ using an inverse for the sum of each factor's bits above position
zero. That sum lies between 0 and 63, so its nonzero field value means the
factor is at least two. These extra rows make the public-input promise robust.

Completeness follows by taking the binary expansions and inverses from every
nontrivial integer factorisation. For soundness, range constraints give
$0\le u,v<2^{64}$, so $uv<2^{128}<P$. Both $uv$ and $n$ are therefore
canonical integers, and the field equality lifts to $uv=n$ over the integers.
The inverse constraints exclude trivial factors. The circuit checks supplied
factors; it does not find them.

\section{Q8: a 4096-bit factorisation}
Put $B=2^{64}$ and represent $u,v,n$ by 64 limbs each:
$u=\sum_i u_iB^i$, $v=\sum_i v_iB^i$, and $n=\sum_i n_iB^i$.
Range-constrain every limb, including public $n_i$, to 64 bits. Introduce 4096
partial-product signals and 129 carries, with rows
\begin{align*}
 u_iv_j&=q_{ij} &&(0\le i,j<64),\\
 \sum_{i+j=k}q_{ij}+c_k&=n_k+Bc_{k+1} &&(0\le k<128),\\
 c_0&=c_{128}=0,
\end{align*}
where $n_k=0$ for $k\ge64$. Constrain $c_1,\ldots,c_{127}$ to 70 bits.
For each factor take the sum of its 4095 bits other than global bit zero and
require an inverse of that sum. Since $4095<P$, this enforces that the factor
is at least two.

\paragraph{Completeness.}
Every integer factorisation $n=uv<2^{4096}$ with $u,v\ge2$ has factors below
$2^{4096}$. Use their base-$B$ limbs, exact partial products, and ordinary
schoolbook carries. Inductively $c_k\le64(B-1)$: with at most 64 products,
the column total is at most
$64(B-1)^2+64(B-1)=64B(B-1)$. Hence $c_{k+1}\le64(B-1)<2^{70}$.
The output fits 64 limbs, so all higher output limbs and the final carry are
zero. Binary decompositions and the indicated inverses complete the witness.

\paragraph{Soundness.}
Every $u_iv_j<B^2=2^{128}<P$, so each partial-product field equation fixes
the exact canonical product $q_{ij}$. The left side of any column is at most
$64(B-1)^2+(2^{70}-1)$; the right side is at most
$(B-1)+B(2^{70}-1)$. Both are nonnegative and less than $2^{135}<P$.
Thus every column equality is an integer equality, not merely a congruence.
Multiplying column $k$ by $B^k$ and summing cancels the intermediate carries,
giving $uv=n+B^{128}c_{128}-c_0=n$. The zero upper outputs prevent overflow
past the 64 public limbs, and the inverse constraints make both factors
nontrivial. This proves both directions independently of witness generation.

\section{Verification record and reproducibility}
The exporter writes only generic quadratic equalities and input declarations.
All private and auxiliary assignments are supplied explicitly. Compilation uses
Circom 2.2.3, BN128, optimisation level zero, and disabled witness sanity checks;
the R1CS remains enforced. Snarkjs 0.7.6 checks valid witnesses. A separate
Python evaluator reads exported compiled rows, and malicious assignments are
evaluated directly against those rows. Witness-generator refusal is never
counted as soundness evidence.

The audit covers range boundaries, nonbinary bits, factor zero and one,
incorrect public products, partial products, carry bits and endpoints, high
limbs, overflow, and modular aliases. An explicit false factorisation
$2((P+7)/2)=7+P$ satisfies multiplication columns with unrestricted field
carries; the released carry ranges reject it. Deliberately removing the relevant
range rows makes this negative control pass, demonstrating why those bounds
are required. Public-limb and factor-range omission controls are checked too.

\subsection{Specified validation inputs}
The 13 Fourier vectors at each size are: the zero vector; the constant
vector $2-3\mathrm{i}$; four impulses of amplitude $1-2\mathrm{i}$ at
positions $0,1,\lfloor N/3\rfloor,N-1$; three modes
$x_j=\exp(2\pi\mathrm{i}kj/N)$ for $k=1,N/4,N-1$; and four complex Gaussian
vectors. For the latter, independent standard normal samples supply the
real and imaginary parts, using the default NumPy random generator with
seeds 101, 202, 303 and 404. The recorded environment uses NumPy 2.4.3 and
Python 3.13. The normalised maximum error is
\[
\varepsilon=\frac{\max_r|y_r-y_r^{\mathrm{ref}}|}
 {\max(1,\max_r|y_r^{\mathrm{ref}}|)},
\]
where the reference is the unnormalised forward transform evaluated by an
independent FFT. A finite error no larger than $10^{-10}$ is required.
Complex arrays use double precision. Small-size checks additionally compare
every basis vector through $N=128$ against the direct exponential formula.
The target runs impose a 900-second limit per vector and a 4 GiB memory
ceiling. Their plaintext elapsed times are validation costs, not estimates
of encrypted execution.

The valid arithmetic witnesses exercise the following inputs, with
$B=2^{64}$. In Q7 and Q8 the public value is the product of the listed factors.
\begin{center}
\begin{tabularx}{\linewidth}{lX}\toprule
Question & Valid inputs\\\midrule
Q5 & $x=0,1,B-1$\\
Q6 & $r=0,2,P-1$\\
Q7 & $(u,v)=(2,3),(3,2),(65537,65539),(2,(B-2)/2)$\\
Q8 & $(u,v)=(2^{2047}+123,2^{2047}+321)$,
$(2,2^{4094}+1)$, $(2^{4094}+1,2)$\\\bottomrule
\end{tabularx}
\end{center}
Valid assignments are checked both before and after compilation. Invalid
assignments are injected directly into the compiled rows, so failure of a
witness-generation routine is not used as rejection evidence. Mutations
exercise nonbinary bits, out-of-range values, wrong inverses, zero and unit
factors, changed public products, partial products, carry bits, endpoint
carries, high factor and public limbs, and overflow beyond 4096 bits.
The 64-bit range interface also rejects noncanonical external integers
$-1,P,P+1$ before they can alias field values.

\subsection{Why the omitted-bound controls are informative}
For the large-factor carry control, take $u=2$, $v=(P+7)/2$ and public
$n=7$. The ordinary product is $P+7$. Write the factors and $n$ in limbs,
set $c_0=0$, and compute successive field carries by
\[
c_{k+1}=B^{-1}\left(\sum_{i+j=k}u_iv_j+c_k-n_k\right)\pmod P.
\]
All multiplication columns are then satisfied. Summing the weighted columns
shows that $c_{128}=B^{-128}(uv-n)=0$ in the field, so even zero endpoints
cannot reject the false integer equality by themselves. The released system
rejects these carries because of its 70-bit range equations. Deleting those
range equations makes this assignment pass, isolating their role.

The small-factor control uses the same modular alias and isolates the
64-bit factor bound. A public-limb control replaces $n_0$ by $n_0+B$ and
$n_1$ by $n_1-1$, leaving the represented integer unchanged but violating
the canonical limb range. A separate overflow control attempts to represent
$2^{4095}\cdot2$ with all public limbs zero. These controls check distinct
ways in which a field assignment could violate the intended integer statement.

\paragraph{Checks behind the operation-network diagrams.}
The query diagram computes the sum and weighted contributions of $(3,8,19)$,
then repeatedly divides the exact answer by 31 and checks recovery of all
three entries. The eight-slot Fourier diagram is generated from three
explicit butterfly matrices and the bit-reversal permutation. Their product
is compared with the direct exponential matrix, and the emitted packed
circuit is checked on all eight basis vectors against that same formula,
with a maximum-error tolerance of $10^{-12}$. The target-size cost bars
and stage groups are read from the recorded search results.
For the decimal multiplication illustration, the successive column totals
are $21,28,10,1$, with output digits $1,8,0,1$ and outgoing carries
$2,2,1,0$. The check verifies that their weighted digit sum equals
$23\cdot47=1081$. These small examples explain the operations; the
large-dimension numerical and compiled checks remain those reported above.

\subsection{Constraint counts}
For Q8 the 192 factor and public limbs contribute $192\cdot65=12480$
rows: 64 bit equations and one reconstruction equation per limb. The
4096 partial products contribute 4096 rows. The 127 intermediate carries
contribute $127\cdot71=9017$ rows: 70 bit equations and one reconstruction
each. Add 128 column equations, two zero-endpoint equations and two
factor lower-bound equations, for $25725$ rows in total. Auxiliary signals
are supplied explicitly. The signal count is
$192\cdot65+4096+129+127\cdot70+2=25597$.
For Q7, three 64-bit decompositions, one product and four inverse equations
give $195+1+4=200$ rows. The 65-row range gadget and one-row exclusion
gadget complete the four families.

\subsection*{Generated results}
\begin{center}\begin{tabular}{rrrrrr}\toprule
$N$ & Multiplies & Adds & Rotations & Depth & Cost (s)\\\midrule
32768 & 4757 & 4754 & 433 & 3 & 29.299516\\
65536 & 13313 & 13310 & 570 & 3 & 78.083040\\
\bottomrule\end{tabular}\end{center}
For $N=32768$, stages are [7, 7, 1], baby sizes [2048, 32, 256], and the dense baseline is 185.962068 seconds (6.347 times the selected cost). The 13 target cases have maximum normalised error $2\times 10^{-15}$. 
For $N=65536$, stages are [8, 7, 1], baby sizes [4096, 32, 256], and the dense baseline is 370.374740 seconds (4.743 times the selected cost). The 13 target cases have maximum normalised error $2.13\times 10^{-15}$. 
\begin{center}\begin{tabular}{rrrrrrrrr}\toprule
$n$ & 1 & 3 & 5 & 7 & 9 & 11 & 13 & 15\\
Gates & 0 & 1 & 7 & 30 & 88 & 205 & 411 & 742\\\bottomrule\end{tabular}\end{center}
\begin{center}\begin{tabular}{lrr}\toprule Circuit & Compiled rows & Signals\\\midrule
Q5 & 65 & 65\\
Q6 & 1 & 2\\
Q7 & 200 & 199\\
Q8 & 25725 & 25597\\
\bottomrule\end{tabular}\end{center}
The compiled audit records 66 passing checks. The challenge and affected Boolean regression suites collect 63 and 525 tests respectively; all pass. The modular enumeration counts for primes 3, 5, 7, 11, 13, 19 are 9, 25, 1, 1, 73, 1.
Tables are generated from the recorded verification results; individual records retain the producing commands and source hashes.
The supplementary index check reconstructs 2730 majority circuit states. The incorrect nine-input circuit disagrees with majority on 54 of 512 states, despite having the same functional inputs and the same number of active indices.


\subsection{Reproduction and scope of the evidence}
Reproduction proceeds in four parts: regenerate the discrete constructions and
compare their complete or symbolic outputs; emit and audit the four quadratic
systems; execute both selected Fourier circuits and compare their results;
and regenerate the tables and documents from the resulting records. The
accompanying reproduction guide supplies executable commands and dependency
pins separately from the mathematical exposition.

The original verification record contains 63 challenge tests and 525 affected
Boolean regression tests. It also records 66 compiled arithmetic checks and
13 plaintext vectors at each target Fourier size. Missing tools, skipped
required tests and unresolved outcomes are failures of full verification.
The new index-deconvolution examples are a separate, bounded manuscript
check, described in Section S4. The original numerical and compiled results
are retained; rebuilding the documents does not rerun those experiments.
Records retain the producing commands, source hashes, numerical tolerances,
input sizes, seeds, elapsed times and observed memory. Document validation
checks successful compilation, resolved references and page bounds. It does
not replace the mathematical or computational checks.

\paragraph{What has not been measured.}
No complete 2025-input circuit has been materialised, no ciphertext Fourier
timing has been taken, and no zero-knowledge proof has been generated.
The arithmetic results establish properties of the compiled constraint
systems and supplied witnesses. They are not a benchmark of proving time,
proof size or a deployed cryptographic protocol.

\begin{thebibliography}{9}
\bibitem{challenge} 0xPARC. \emph{Beyond Primitive Computing: Puzzles and Connections}, Field Extension.
\url{https://fieldextension.org/3b1b/}.
\bibitem{cohen} G. Cohen, I. B. Damg\aa rd, Y. Ishai, J. K\"olker,
P. B. Miltersen, R. Raz, and R. D. Rothblum.
\emph{Efficient Multiparty Protocols via Log-Depth Threshold Formulae}, 2013.
\url{https://www.wisdom.weizmann.ac.il/~/ranraz/publications/Pmajority_mpc-1.pdf}.
\bibitem{cheon} J. H. Cheon, K. Han, and M. Hhan.
\emph{Faster Homomorphic Discrete Fourier Transforms and Improved FHE
Bootstrapping}. Cryptology ePrint 2018/1073.
\url{https://eprint.iacr.org/2018/1073}.
\bibitem{circom} iden3. \emph{Circom 2.2.3} and \emph{snarkjs 0.7.6}.
\url{https://docs.circom.io/};
\url{https://www.npmjs.com/package/snarkjs/v/0.7.6}.
\end{thebibliography}
\end{document}
